\documentclass[main.tex]{subfiles}
\begin{document}
\section{Semantics}

We define 2 languages in this section, \surface{} and
\core. \surface{} corresponds to the \texttt{SkCanvas} API a user of Skia
typically uses to draw pictures.

\todo[inline]{As this is a description of the Skia API, maybe Draw needs to be
  specific as in DrawRect etc. and not the generic Geometry Draw}

\begin{center}
  \begin{bnf}[
    colspec = {llcll},
    column{1} = {font = \sffamily},
    column{2} = {mode = dmath},
    column{4} = {font = \ttfamily},
    ]
    s : \surface ::=
    | $c$
    ;;
    c : Command  ::=
    | Draw($g$, $p$)
    | Clip($cl$, Intersect | Difference)
    | Save() ; $c$ ; Restore()
    | SaveLayer($p$) ; $c$ ; Restore()
    | $c$ ; $c$
    ;;
    g : DrawGeometry ::=
    | Rect($\ldots$) | RRect ($\ldots$)
    | Path($\ldots$) | $\ldots$
    ;;
    cl : ClipGeometry ::=
    | Rect($\ldots$) | RRect ($\ldots$)
    | Path($\ldots$)
    ;;
    p : Paint ::=
    | Paint($f$, $st$, $c_f$, $b$)
    ;;
    f : Fill ::=
    | Color(Pixel)
    | LinearGradient($\ldots$)
    | $\ldots$
    ;;
    st : Style ::=
    | Solid | Stroke
    ;;
    c_f : Filter ::=
    | IdFilter | LumaFilter | $\ldots$
    ;;
    b : BlendMode ::=
    | SrcOver | DstIn | Src | $\ldots$
  \end{bnf}
\end{center}

A \surface{} program maintains 2 types of state, the state of canvas and the clip
state. The canvas state describes what is currently drawn on the drawing
surface. The clip state constrains the extent of what is drawn on the drawing
surface. Consider the program:

\begin{verbatim}
Draw(Rect(64, 64, 256, 256), RED_PAINT)
\end{verbatim}

Here we are drawing on a canvas of size 256$\times$256 and \texttt{RED\_PAINT}
being a \texttt{Paint} that just fills a geometry with the color {\color{red}
  Red}. This program would render to the following image:


\begin{center}
\frame{\includegraphics[scale=0.3]{drawing1.pdf}}
\end{center}

Now consider the program:

\begin{verbatim}
Clip(Rect(0, 0, 192, 192));
Draw(Rect(64, 64, 256, 256), RED_PAINT)
\end{verbatim}

Here we clip to the upper-right corner of the square. The resulting image for
rendering this program is:

\begin{center}
\frame{\includegraphics[scale=0.3]{drawing2.pdf}}
\end{center}

The contents of \texttt{Paint} merit further discussion. \texttt{Fill}s describe
what pixels to fill the geometry with, which could be the same color or a shader
describing say a gradient. \texttt{Style} describes how the geometry is drawn,
whether the extent of the geometry is filled, or only its bounds.
\texttt{Filter}s basically are functions to transform pixels. Finally
\texttt{BlendMode}s describe how the drawn and transformed geometry is
\textit{composited} or \textit{blended} onto the canvas. You can consider a
blend mode as an operator over pixels.

\todo[inline]{I ran out of gas here. Probably want pictures for all of this? We
  cannot assume readers have ever used Skia before, so they do not know what I
  am on about in the previous paragraph.}

It is important to undertsnad that \texttt{Clip}s are additive to current clip
state of the canvas. Just with \texttt{Clip}s alone you may not reset the clip
state. This is taken care of by the \texttt{Save(); ...; Restore()} and
\texttt{SaveLayer(); ...; Restore()} commands. Calling a
\texttt{Save}/\texttt{SaveLayer} is like saving a checkpoint of the current clip
state. When a matching \texttt{Restore()} is called, all changes made to the clip
state are discarded. Internally Skia implements this checkpointing system by
managing a stack of clip states, with \texttt{Save}/\texttt{SaveLayer} pushsing
a copy of the current clip state onto the stack, while a \texttt{Restore} pops
the latest clip state of the stack.

While \texttt{Save}/\texttt{SaveLayer} both manipulate the clip state of the
canvas, \texttt{SaveLayer} also allocates a new seperate surface. All commands
between a \texttt{SaveLayer} and a matching \texttt{Restore} affect the new
surface. Once \texttt{Restore} is called, the new surface is merged onto the old
one. How this merge happens is dictated by the \texttt{BlendMode} passed to the
\texttt{Paint} of the \texttt{SaveLayer}. This functions just like \texttt{Draw}
commands. \texttt{SaveLayer}s can also be used to modify the entire new surface
before it is blended back onto the old surface. The \texttt{Paint} of a
\texttt{SaveLayer} can be passed an alpha value, which is multiplied over the
entire new surface right before blending. In a similiar vein, a \texttt{Filter}
inside the \texttt{Paint} will be applied over the entire surface before
blending. It is important to note that complicated fills like gradients, and
fill/stroke styles do not affect how a \texttt{SaveLayer} blends its surface
onto the previous one.

\todo[inline]{I dont like old/new/previous here. We always talk about surfaces
  as bottom and top say for example in the lean formalization. Also probably
  need to change surfaces to layers here, or layers to surfaces everywhere
  else.}

As our main goal is to try and verify rewrites in Skia, it is hard to do so over
\surface{}. This is mainly because of the hidden surface and clip state of
\surface{}. That is why we introduce a new language called \core{} which is a
declrative version of \surface{}.\;\core{} makes surfaces the main value type of
the language, and also make explicit the clip state at each draw command. This
means neither \texttt{Save}/\texttt{SaveLayer} are needed to ``manage'' clip
state in \core{}. This also has the consequence of completely removing
\texttt{Save}s from \core{}. \texttt{SaveLayers} remain in \core{} because they
also allow us to combine surfaces in addition to managing the clip state.

\todo[inline]{Bad paragraph. We still don't know if this paper is verification
  first or optimization first. Also there are way to many \surface{} and \core{}
  in this paragraph.}

\begin{center}
  \begin{bnf}[
    colspec = {llcll},
    column{1} = {font = \sffamily},
    column{2} = {mode = dmath},
    column{4} = {font = \ttfamily},
    ]
    s : \core ::=
    | $l$
    ;;
    l : Layer ::=
    | Empty()
    | Draw($l$, $g$, $p$, $cl$)
    | SaveLayer($l$, $l$, $p$)
    ;;
    g : Geometry ::=
    | Rect($\ldots$) | RRect ($\ldots$)
    | Path($\ldots$) | $\ldots$
    ;;
    cl : ClipGeometry ::=
    | Rect($\ldots$) | RRect ($\ldots$)
    | Path($\ldots$) | Intersect($cl$, $cl$)
    ;;
    p : Paint ::=
    | Paint($f$, $st$, $c_f$, $b$)
    ;;
    f : Fill ::=
    | Color(Pixel)
    | LinearGradient($\ldots$)
    | $\ldots$
    ;;
    st : Style ::=
    | Solid | Stroke
    ;;
    c_f : Filter ::=
    | IdFilter | LumaFilter | $\ldots$
    ;;
    b : BlendMode ::=
    | SrcOver | DstIn | Src | $\ldots$
  \end{bnf}
\end{center}

The main syntactical structures of \core{} are \textsf{Layer}s. A \texttt{Layer}
correpsonds to one state of the draw surface. An \textsf{Empty} layer is a draw
surface with all the pixels set to the argb color \texttt{(0, 0, 0, 0)}. A
\textsf{Draw} layer draws a geometry onto the given layer. Just like the
\textsf{Draw*} commands in \surface, how the shape is drawn and blended onto the
layer is governed by the \textsf{Paint}. A \textsf{SaveLayer} layer is similiar
to a \textsf{SaveLayer(); ...; Restore();} pair in \surface{}.

\begin{minipage}{0.48\textwidth}
\begin{verbatim}
bottom
...
SaveLayer(paint);
top
...
Restore();
\end{verbatim}
\end{minipage}\hfill
\begin{minipage}{0.48\textwidth}
\begin{verbatim}
SaveLayer(bottom, top, paint)
\end{verbatim}
\end{minipage}

\textsf{Paint}s behave the same way in \core{} as they do in \surface{}. Now
that we have syntax for \core{} we can define a denotational semantics over it.

\todo[inline]{everything after this i sprobably written in the worst way
  possible}

We define a denotation function over \texttt{Layer}s called
\denote{\;\cdot\;}{Layer}. Its signature is:

\begin{equation*}
  \denote{\;\cdot\;}{Layer} : int^2 \to int^4
\end{equation*}

That is the denotation of Layer is map from 2D points to ARGB pixels. We now
define the denotations of the various constructors of \texttt{Layer}.

\begin{equation*}
  \denote{\;\texttt{Empty()}\;}{Layer} = \lambda pt . (0, 0, 0, 0)
\end{equation*}

\begin{align*}
  \denote{\;\texttt{Draw(l, g, (f, sty, cf, bm), cl)}\;}{Layer} = \lambda pt .
  bm (l pt) (cf (if (style g) and (clip pt) then (fill pt) else (0, 0, 0, 0)))
\end{align*}

\todo[inline]{I have not added savelayer because it will be to painful to looka
  t at as is}

\begin{equation*}
  \denote{\;\texttt{Empty()}\;}{Layer} = \lambda pt . (0, 0, 0, 0)
\end{equation*}






\end{document}
